\section{Agenda de investigación: tres frentes de avance para la siguiente generación de Agents multimodales}

\subsection{Frente de avance uno: capa de representación unificada}
El curso discutió extensamente las capturas de pantalla, los árboles de accesibilidad, el código de acciones y las cadenas de razonamiento, que en esencia son representaciones intermedias de distintas modalidades. Una dirección importante para la siguiente etapa es construir una capa de representación unificada que reduzca la pérdida de información al cambiar de modalidad.

\begin{knowledgebox}{Problemas que debe resolver la capa de representación unificada}
\begin{itemize}
    \item Cómo mapear los objetos visuales, los objetos semánticos y los objetos de acción a una misma estructura;
    \item Cómo lograr que distintas plataformas (web, escritorio, móvil) compartan una interfaz de abstracción;
    \item Cómo dar soporte a la compresión de estado y la recuperación reversible en tareas de cadena larga.
\end{itemize}
\end{knowledgebox}

\subsection{Frente de avance dos: planificación verificable de largo alcance}
Los modelos actuales todavía caen fácilmente en estrategias miopes ante tareas complejas. La parte final del curso enfatizó la contribución del planner, lo que indica que la «planificación verificable» será el próximo foco de competencia: no basta con generar un plan, sino verificar de manera continua si el plan sigue siendo coherente con el estado del entorno.

\begin{importantbox}{La evaluación de la capacidad de planificación debe pasar de «tiene o no tiene» a «calidad»}
Las evaluaciones futuras no solo deben preguntar si el modelo puede producir un plan, sino también evaluar la viabilidad de ejecución del plan, su capacidad de actualizarse y su eficiencia de recuperación ante errores. Sin estos indicadores, el planner puede reducirse a un texto intermedio meramente decorativo.
\end{importantbox}

\subsection{Frente de avance tres: seguridad y gobernanza nativas}
Cuando un Agent comienza a operar cuentas reales y sistemas de negocio, la seguridad ya no puede posponerse. Los límites de permisos, los registros de auditoría, la reversión ante anomalías y los mecanismos de intervención humana deben convertirse en restricciones nativas del entrenamiento y el despliegue del modelo, y no en parches previos al lanzamiento.

\begin{warningbox}{El mayor riesgo de un Agent no es solo «no poder hacerlo», sino «hacerlo mal y sin poder rastrearlo»}
Si un sistema carece de una cadena auditable, después de un incidente resulta imposible reconstruir los límites de responsabilidad. Aunque el curso se centra en la investigación, su pensamiento sobre entornos y evaluación ya ofrece una base práctica para la gobernanza nativa.
\end{warningbox}

\subsection{Lista de tareas ejecutables para los próximos 12 meses}
\begin{table}[H]
\centering
\begin{tabular}{p{2.8cm}p{4.8cm}p{3.6cm}}
\toprule
\textbf{Dirección} & \textbf{Tarea propuesta} & \textbf{Resultado esperado} \\
\midrule
Representación unificada & Construir una capa intermedia de objeto-acción entre plataformas & Reducir el costo de migración y la tasa de error en las acciones \\
Optimización de la planificación & Introducir verificación en línea del plan y mecanismos de replanificación & Aumentar la tasa de finalización de tareas largas \\
Volante de datos & Agrupar automáticamente las muestras de fallo por categorías y reincorporarlas al entrenamiento & Reducir la tasa de recurrencia de fallos frecuentes \\
Gobernanza de seguridad & Establecer plantillas de permisos, auditoría y estrategias de reversión & Sostener un despliegue en producción controlable \\
\bottomrule
\end{tabular}
\caption{Hoja de ruta que lleva la metodología del curso a la capa de ejecución del equipo}
\end{table}

\subsection{Resumen de la sección}
La competencia de la siguiente generación de Agents multimodales no dependerá solo de la escala de parámetros del modelo, sino de quién logre convertir la representación unificada, la planificación de largo alcance y los mecanismos de gobernanza en capacidades de ingeniería reutilizables.

